% LTeX: language=de

\subsubsection{Aufgaben}

\begin{Exercise}{S. 32/6}{}
\begin{enumerate}[label=\arabic*.,leftmargin=*]
    \item [6.] {
        $f(x) = \dfrac{1}{x^{2} + 1}$; \qquad $g(x) = \arctan\left(f(x)\right)$; \qquad $D_{f} = D_{g} = \mathbb{R}$ \\[-1.0ex] 
        \begin{enumerate}[label=\alph*),leftmargin=*]
        \item Zeichnung: \\[-3.0ex]{
            \raisebox{-\height}{
            \begin{tikzpicture}[
                    declare function={
                        f_f(\x) = (1/((\x*\xScale)^(2) + 1)))*(1/\yScale);
                    }
                ]

                % Set variables
                \pgfmathsetmacro{\axisPosWidth}{6.0};
                \pgfmathsetmacro{\axisNegWidth}{-6.0};
                \pgfmathsetmacro{\axisPosHeight}{3.0};
                \pgfmathsetmacro{\axisNegHeight}{-1.0};
                \pgfmathsetmacro{\xIncrement}{1.0};
                \pgfmathsetmacro{\yIncrement}{1.0};
                \pgfmathsetmacro{\xScale}{1.0};
                \pgfmathsetmacro{\yScale}{0.5};
                \pgfmathsetmacro{\gridxIncrement}{0.5};
                \pgfmathsetmacro{\gridyIncrement}{0.5};

                % Axis/grid limits
                \pgfmathsetmacro{\xUpperLimitAxis}{\xIncrement * floor(\axisPosWidth / \xIncrement)}
                \pgfmathsetmacro{\xLowerLimitAxis}{\xIncrement * ceil(\axisNegWidth / \xIncrement)}
                \pgfmathsetmacro{\yUpperLimitAxis}{\yIncrement * floor(\axisPosHeight /\yIncrement)}
                \pgfmathsetmacro{\yLowerLimitAxis}{\yIncrement * ceil(\axisNegHeight / \yIncrement)}

                \pgfmathsetmacro{\xUpperLimitGrid}{\gridxIncrement * floor(\axisPosWidth / \gridxIncrement)}
                \pgfmathsetmacro{\xLowerLimitGrid}{\gridxIncrement * ceil(\axisNegWidth / \gridxIncrement)}
                \pgfmathsetmacro{\yUpperLimitGrid}{\gridyIncrement * floor(\axisPosHeight / \gridyIncrement)}
                \pgfmathsetmacro{\yLowerLimitGrid}{\gridyIncrement * ceil(\axisNegHeight / \gridyIncrement)}

                % Axis environment (PGFPlots)
                \begin{axis}[
                    xmin=\xLowerLimitGrid, xmax=\xUpperLimitGrid,
                    ymin=\yLowerLimitGrid, ymax=\yUpperLimitGrid,
                    axis lines=middle,
                    axis line style={->, >=stealth},
                    grid=both,
                    xlabel={$x$},
                    ylabel={$y$},
                    xlabel style={anchor=west, at={(current axis.right of origin)}},
                    ylabel style={anchor=south, at={(current axis.above origin)}},
                    xtick={\xLowerLimitAxis,\xLowerLimitAxis+1,...,\xUpperLimitAxis},
                    ytick={\yLowerLimitAxis,\yLowerLimitAxis+1,...,\yUpperLimitAxis},
                    xticklabel={\pgfmathparse{\tick*\xScale}\pgfmathprintnumber{\pgfmathresult}},
                    yticklabel={\pgfmathparse{\tick*\yScale}\pgfmathprintnumber{\pgfmathresult}},
                    tick label style={font=\scriptsize},
                    x=1cm,
                    y=1cm,
                    minor tick num=1,
                    scale only axis,
                    grid=both,
                    restrict y to domain=\yLowerLimitGrid:\yUpperLimitGrid,
                ]
                    
                % Graph limits
                \pgfmathsetmacro{\xLowerLimitGraph}{\xLowerLimitAxis}
                \pgfmathsetmacro{\xUpperLimitGraph}{\xUpperLimitAxis}
                
                % G_f
                \pgfmathsetmacro{\Gfax}{0.5*(1/\xScale)}
                \pgfmathsetmacro{\Gfay}{f_f(\Gfax)}
                \addplot[thick, myb, samples=1000, domain=\xLowerLimitGraph:\xUpperLimitGraph]
                    {f_f(x)};
                \node[anchor=south west] at (axis cs:{\Gfax},{\Gfay}) {\textcolor{myb}{$G_{f}$}};

                % h_1 für y1
                \pgfmathsetmacro{\Gtby}{0*(1/\yScale)}
                \addplot[thick, myorange] coordinates {(\xLowerLimitGraph,\Gtby) (\xUpperLimitGraph,\Gtby)};
                \node[anchor=north west] at (axis cs:{0},{\Gtby}) {\textcolor{myorange}{$G_{h_{1}}$}};

                \end{axis}
            \end{tikzpicture}
            }\\[1.0ex]
            \begin{itemize}[leftmargin=*]
                \item Symmetrie: {
                    \begin{minipage}[t]{\linewidth}
                        $D_{f}$ ist achsensymmetrisch zur y-Achse, da $f(x) = f(-x)$ gilt\\[-2.0ex]
                    \end{minipage}
                    }
                \item Globalverlauf: \\[-2.0ex]{
                    \begin{align*}
                        \lim\limits_{x \to -\infty} \arctan\left(\overbrace{f(x)}^{\textcolor{myg}{0^{+}}}\right) &= 0^{+}; \qquad \qquad \lim\limits_{x \to +\infty} \arctan\left(\overbrace{f(x)}^{\textcolor{myg}{0^{+}}}\right) &= 0^{+} &&
                    \end{align*}
                    }
                \item Asymptoten: \\[-3.0ex]{
                    \begin{itemize}
                        \item Waagerechte Asymptoten: \quad {
                            \vspace{-5.0ex}
                            \begin{align*}
                                h_{1}(x): \quad &y = 0
                            \end{align*}
                            }
                    \end{itemize}
                    }
                \item Schnittpunkte mit den Koordinatenachsen: \\[-2.0ex]{
                    \begin{itemize}
                        \item $x$-Achse: \quad keine, da $\arctan\left(0\right) = 0$ für $f(x) = 0$, aber $f(x) \neq 0$\\[-1.0ex]
                        \item $y$-Achse: \quad $g(0) = \dfrac{\pi}{4}$\\[-1.5ex]
                    \end{itemize}
                    }
                \item Extrempunkte: \\[1.0ex]{
                    \texthl{Arcustangensfunktion ist \textbf{sms}, daher überträgt sich die Montonie von $G_{f}$ auf $G_{g}$}\\[1.0ex]
                    $\Rightarrow$ \quad $G_{g}$ ist sms in $\left] \ -\infty; 1 \ \right]$ \\[1.0ex]
                    $\Rightarrow$ \quad $G_{g}$ ist smf in $\left[ \ 1; \infty \ \right[$ \\[-2.0ex]
                    \begin{itemize}
                        \item Hochpunkte: \quad {
                            $H_{1}\left(0 / \frac{\pi}{4}\right)$
                            }
                        \item Tiefpunkte: \quad {
                            keine
                        }
                        \item Randextrema: \quad {
                            keine
                        }
                    \end{itemize}
                    }
                \item Wendepunkte: \quad {
                    $W_{1} \left( -\SI{0.50}{} \ / \ g\left(-\SI{0.50}{}\right)\right)$; \qquad  $W_{2} \left( \SI{0.50}{} \ / \ g\left(\SI{0.50}{}\right)\right)$
                }
            \end{itemize}
            }
        \end{enumerate}
        }
\end{enumerate}
\end{Exercise}

\begin{Exercise}{S. 107/2}{}
\begin{enumerate}[label=\arabic*.,leftmargin=*]
    \item [2.] {
        $g(x) =  \sqrt[2]{f(x)}$ \\[-3.0ex] 
        \begin{itemize}[leftmargin=*]
            \item Definitionsmenge: \\[1.0ex]{
                $g(x)$ definiert für $f(x) \geq 0$\\[1.0ex]
                $\Rightarrow$ \quad $D_{g} = \left] \ -\infty; -3 \ \right]$ \quad $\cup$ \quad $\left[ \ -1; 0 \ \right[$ \quad $\cup$ \quad $\left] \ 0; +\infty \ \right[$\\[-2.0ex]
                }
            \item Verhalten am Rand der Definitionsmenge: \\[-2.0ex]{
                \begin{align*}
                    \lim\limits_{x \to -\infty} \sqrt[2]{\underbrace{f(x)}_{\textcolor{myg}{0^{+}}}} &= 0^{+}; \qquad \qquad &\lim\limits_{x \to +\infty} \sqrt[2]{\underbrace{f(x)}_{\textcolor{myg}{0^{+}}}} &= 0^{+} && \\[1.5ex]
                    g(-3) &= 0; &g(-1) &= 0 && \\[2.0ex]
                    \lim\limits_{x \to 0^{-}} \sqrt[2]{\underbrace{f(x)}_{\textcolor{myg}{\infty}}} &= + \infty; \qquad \qquad &\lim\limits_{x \to 0^{+}} \sqrt[2]{\underbrace{f(x)}_{\textcolor{myg}{\infty}}} &= + \infty && \\[-2.0ex]
                \end{align*}
                }
            \item Nullstellen: \\[1.0ex]{
                $g(x) = 0$ für $f(x) = 0$ \qquad $\Rightarrow$ \quad $x_{1} = -3$; \quad $x_{2} = -1$\\[-2.0ex]
                }
            \item Extremstellen: \\[1.0ex]{
                Wurzelfunktion ist \textbf{sms}, daher überträgt sich die Montonie von $G_{f}$ auf $G_{g}$\\[1.0ex]
                $\Rightarrow$ \quad $G_{g}$ ist sms in $\left] \ -\infty; x_{H_{1}} \ \right]$ sowie $\left[ \ -1; 0 \ \right[$ \quad, wobei gilt $-\infty < x_{H_{1}} < -3$ \\[1.0ex]
                $\Rightarrow$ \quad $G_{g}$ ist smf in $\left[ \ x_{H_{1}} ; -3 \ \right]$ sowie $\left] \ 0; \infty \ \right[$ \\[-2.0ex]
                \begin{itemize}
                    \item Hochstellen: \quad {
                        nicht eindeutig bestimmbar aus $G_f$
                        }
                    \item Tiefstellen: \quad {
                        $x_{1} = -3$; \quad $x_{2} = -1$
                    }\\
                \end{itemize}
                }
        \end{itemize}
        $h(x) =  \ln\left(f(x)\right)$ \\[-3.0ex] 
        \begin{itemize}[leftmargin=*]
            \item Definitionsmenge: \\[1.0ex]{
                $h(x)$ definiert für $f(x) > 0$\\[1.0ex]
                $\Rightarrow$ \quad $D_{h} = \left] \ -\infty; -3 \ \right[$ \quad $\cup$ \quad $\left] \ -1; 0 \ \right[$ \quad $\cup$ \quad $\left] \ 0; +\infty \ \right[$\\[-2.0ex]
                }
            \item Verhalten am Rand der Definitionsmenge: \\[-2.0ex]{
                \begin{align*}
                    \lim\limits_{x \to -\infty} \ln\left(\underbrace{f(x)}_{\textcolor{myg}{0^{+}}}\right) &= - \infty; \qquad \qquad &\lim\limits_{x \to +\infty} \ln\left(\underbrace{f(x)}_{\textcolor{myg}{0^{+}}}\right) &= - \infty  && \\[1.5ex]
                    \lim\limits_{x \to -3^{-}} \ln\left(\underbrace{f(x)}_{\textcolor{myg}{0^{+}}}\right) &= - \infty; \qquad \qquad &\lim\limits_{x \to -1^{+}} \ln\left(\underbrace{f(x)}_{\textcolor{myg}{0^{+}}}\right) &= - \infty && \\[1.5ex]
                    \lim\limits_{x \to 0^{-}} \ln\left(\underbrace{f(x)}_{\textcolor{myg}{+\infty}}\right) &= + \infty; \qquad \qquad &\lim\limits_{x \to 0^{+}} \ln\left(\underbrace{f(x)}_{\textcolor{myg}{+\infty}}\right) &= + \infty && \\[-2.0ex]
                \end{align*}
                }
            \item Nullstellen: \\[1.0ex]{
                $h(x) = 0$ für $f(x) = 1$ \qquad $\Rightarrow$ \quad $x_{1} = \SI{-0.95}{}$; \quad $x_{2} = 1$\\[-2.0ex]
                }
            \item Extremstellen: \\[1.0ex]{
                Logarithmusfunktion ist \textbf{sms}, daher überträgt sich die Montonie von $G_{f}$ auf $G_{h}$\\[1.0ex]
                $\Rightarrow$ \quad $G_{h}$ ist sms in $\left] \ -\infty; x_{H_{1}} \ \right]$ sowie $\left] \ -1; 0 \ \right[$ \quad, wobei gilt $-\infty < x_{H_{1}} < -3$ \\[1.0ex]
                $\Rightarrow$ \quad $G_{h}$ ist smf in $\left[ \ x_{H_{1}} ; -3 \ \right[$ sowie $\left] \ 0; \infty \ \right[$ \\[-2.0ex]
                \begin{itemize}
                    \item Hochstellen: \quad {
                        nicht eindeutig bestimmbar aus $G_f$
                        }
                    \item Tiefstellen: \quad {
                        keine vorhanden
                    }\\
                \end{itemize}
                }
        \end{itemize}
        $i(x) =  e^{f(x)}$ \\[-3.0ex] 
        \begin{itemize}[leftmargin=*]
            \item Definitionsmenge: \\[1.0ex]{
                $i(x)$ definiert für $f(x)$\\[1.0ex]
                $\Rightarrow$ \quad $D_{i} = \mathbb{R} \setminus \{0\}$\\[-2.0ex]
                }
            \item Verhalten am Rand der Definitionsmenge: \\[-2.0ex]{
                \begin{align*}
                    \lim\limits_{x \to -\infty} e^{\overbrace{f(x)}^{\textcolor{myg}{0^{+}}}} &= 1^{+}; \qquad \qquad &\lim\limits_{x \to +\infty} e^{\overbrace{f(x)}^{\textcolor{myg}{0^{+}}}} &= 1^{+} && \\[1.5ex]
                    \lim\limits_{x \to 0^{-}} e^{\overbrace{f(x)}^{\textcolor{myg}{+\infty}}} &= + \infty; \qquad \qquad &\lim\limits_{x \to 0^{+}} e^{\overbrace{f(x)}^{\textcolor{myg}{+\infty}}} &= + \infty && \\[-2.0ex]
                \end{align*}
                }
            \item Nullstellen: \\[1.0ex]{
                $i(x) = 0$ für kein $f(x)$ \qquad $\Rightarrow$ keine Nullstellen\\[-2.0ex]
                }
            \item Extremstellen: \\[1.0ex]{
                e-Funktion ist \textbf{sms}, daher überträgt sich die Montonie von $G_{f}$ auf $G_{i}$\\[1.0ex]
                $\Rightarrow$ \quad $G_{i}$ ist sms in $\left] \ -\infty; x_{H_{1}} \ \right]$ sowie $\left[ \ -\SI{1.30}{}; 0 \ \right[$ \quad, wobei gilt $-\infty < x_{H_{1}} < -3$ \\[1.0ex]
                $\Rightarrow$ \quad $G_{i}$ ist smf in $\left[ \ x_{H_{1}} ; -\SI{1.30}{} \ \right]$ sowie $\left] \ 0; \infty \ \right[$ \\[-2.0ex]
                \begin{itemize}
                    \item Hochstellen: \quad {
                        nicht eindeutig bestimmbar aus $G_f$
                        }
                    \item Tiefstellen: \quad {
                        $x_{1} = -\SI{1.30}{}$
                    }\\
                \end{itemize}
                }
        \end{itemize}
        $j(x) =  \arctan\left(f(x)\right)$ \\[-3.0ex] 
        \begin{itemize}[leftmargin=*]
            \item Definitionsmenge: \\[1.0ex]{
                $j(x)$ definiert für $f(x)$\\[1.0ex]
                $\Rightarrow$ \quad $D_{j} = \mathbb{R} \setminus \{0\}$\\[-2.0ex]
                }
            \item Verhalten am Rand der Definitionsmenge: \\[-2.0ex]{
                \begin{align*}
                    \lim\limits_{x \to -\infty} \arctan\left(\overbrace{f(x)}^{\textcolor{myg}{0^{+}}}\right) &= 0^{+}; \qquad \qquad &\lim\limits_{x \to +\infty} \arctan\left(\overbrace{f(x)}^{\textcolor{myg}{0^{+}}}\right) &= 0^{+} && \\[1.5ex]
                    \lim\limits_{x \to 0^{-}} \arctan\left(\overbrace{f(x)}^{\textcolor{myg}{+ \infty}} \right) &= + {\dfrac{\pi}{2}}^{+}; \qquad \qquad &\lim\limits_{x \to 0^{+}} \arctan\left(\overbrace{f(x)}^{\textcolor{myg}{+ \infty}} \right) &= + {\dfrac{\pi}{2}}^{+} && \\[-2.0ex]
                \end{align*}
                }
            \item Nullstellen: \\[1.0ex]{
                $j(x) = 0$ für $f(x) = 0$ \qquad $\Rightarrow$ \quad $x_{1} = -3$; \quad $x_{2} = -1$\\[-2.0ex]
                }
            \item Extremstellen: \\[1.0ex]{
                Arcustangensfunktion ist \textbf{sms}, daher überträgt sich die Montonie von $G_{f}$ auf $G_{j}$\\[1.0ex]
                $\Rightarrow$ \quad $G_{h}$ ist sms in $\left] \ -\infty; x_{H_{1}} \ \right]$ sowie $\left[ \ -\SI{1.30}{}; 0 \ \right[$ \quad, wobei gilt $-\infty < x_{H_{1}} < -3$ \\[1.0ex]
                $\Rightarrow$ \quad $G_{h}$ ist smf in $\left[ \ x_{H_{1}} ; -\SI{1.30}{} \ \right]$ sowie $\left] \ 0; \infty \ \right[$ \\[-2.0ex]
                \begin{itemize}
                    \item Hochstellen: \quad {
                        nicht eindeutig bestimmbar aus $G_f$
                        }
                    \item Tiefstellen: \quad {
                        $x_{1} = -\SI{1.30}{}$
                    }
                \end{itemize}
                }
        \end{itemize}
        }
\end{enumerate}
\end{Exercise}

\begin{Exercise}{S. 108/4}{}
\begin{enumerate}
    \item [4.]\begin{enumerate}[label=\alph*),leftmargin=*]
        \item[g)] {
            $f(x) = \arctan\left(e^{x}\right)$; \qquad $u(x) = e^{x}$ \\[-2.0ex]
            \begin{itemize}[leftmargin=*]
                \item Definitionsmenge: \\[1.0ex]{
                    $f(x)$ definiert für $u(x)$\\[1.0ex]
                    $\Rightarrow \quad \uuline{D_{f} = \mathbb{R}}$;
                    }
                \item Symmetrie: \\[-2.0ex]{
                    \begin{itemize}
                        \item Achsensymmetrie zur y-Achse: \quad $f(-x) = f(x)$\\[-2.0ex]{
                            \begin{align*}
                                \arctan\left(e^{-x}\right) &= \arctan\left(e^{x}\right) \quad \text{\textit{(falsch)}} &&
                            \end{align*}
                            }
                        \item Punktsymmetrie zum Ursprung: \quad $f(-x) = -f(x)$\\[-2.0ex]{
                            \begin{align*}
                                \arctan\left(e^{-x}\right) &= -\arctan\left(e^{x}\right) && \\[1.5ex]
                                \arctan\left(e^{-x}\right) &= \arctan\left(-e^{x}\right) \quad \text{\textit{(falsch)}} &&
                            \end{align*}
                            }
                    \end{itemize}
                    $\Rightarrow$ \quad keine Symmetrie\\[-2.0ex]
                    }
                \item Grenzverhalten: \\[-3.0ex]{
                    \begin{align*}
                        \lim\limits_{x \to -\infty} \arctan\left(\overbrace{e^{x}}^{\textcolor{myg}{0^{+}}}\right) &= 0^{+}; \qquad \qquad \lim\limits_{x \to +\infty} \arctan\left(\overbrace{e^{x}}^{\textcolor{myg}{+\infty}}\right) &= + {\dfrac{\pi}{2}}^{-} &&\\[-3.0ex]
                    \end{align*}
                    }
                \item Asymptoten: \\[-3.0ex]{
                    \begin{itemize}
                        \item Waagerechte Asymptoten: \quad {
                            \vspace{-5.0ex}
                            \begin{align*}
                                h_{1}(x): \quad &y = 0\\[1.0ex]
                                h_{2}(x): \quad &y = \dfrac{\pi}{2}
                            \end{align*}
                            }
                    \end{itemize}
                    }
                \item Schnittpunkte mit den Koordinatenachsen: \\[-2.0ex]{
                    \begin{itemize}
                        \item $y$-Achse: \quad $f(0) = \arctan\left(e^{0}\right) = \dfrac{\pi}{4}$; \qquad $P_{1} \left( 0 \ / \ \dfrac{\pi}{4} \right)$\\[-1.0ex]
                        \item $x$-Achse: \quad $f(x) = 0$ \\[-2.0ex]{
                            \begin{align*}
                                \arctan\left(e^{x}\right) &= 0 &&\\[1.5ex]
                                e^{x} &= 0 &&\\[1.5ex]
                                x&= \ln\left(0\right)\\[1.5ex]
                                & \textcolor{myr}{\text{n. d.}}\\[-1.5ex]
                            \end{align*}
                            }
                    \end{itemize}
                    }
                \item Extrempunkte: \\[-3.0ex]{
                    \begin{align*}
                        f'(x) &= \dfrac{1}{1 + \left(e^{x}\right)^{2}} \cdot e^{x} \cdot 1&&\\[1.5ex]
                        &= \dfrac{1}{1 + e^{2x}} \cdot e^{x} &&\\[1.5ex]
                        &= \dfrac{e^{x}}{e^{2x} + 1} &&\\[-2.0ex]
                    \end{align*}
                    $z(x) = e^{x}$; \qquad $n(x) = e^{2x} + 1$\\[1.0ex]
                    $f'(x) = 0 \qquad \Rightarrow \qquad z(x) > 0 \qquad \Rightarrow$ \quad keine Nullstellen\\[1.0ex]
                    $n(x) > 0 \qquad \Rightarrow \qquad f'(x) > 0$\\[1.0ex]
                    $\Rightarrow \quad G_{f}$ ist sms in $\mathbb{R}$\\[1.0ex]
                    $\Rightarrow \quad G_{f}$ hat keine Extrempunkte\\[-2.0ex]
                    }
                \item Wendepunkte: \\[-2.0ex]{
                    \begin{align*}
                        f''(x) &= \dfrac{\left(e^{2x} + 1\right) \cdot e^{x} \cdot 1 - e^{x} \cdot e^{2x} \cdot 2}{\left(e^{2x} + 1\right)^{2}}&&\\[1.5ex]
                        &= \dfrac{e^{x} \cdot \left[\left(e^{2x} + 1\right) - 2e^{2x}\right]}{\left(e^{2x} + 1\right)^{2}}&&\\[1.5ex]
                        &= \dfrac{e^{x} \cdot \left(- e^{2x} + 1\right)}{\left(e^{2x} + 1\right)^{2}}&&\\[-1.5ex]
                    \end{align*}
                    $z(x) = e^{x} \cdot \left(- e^{2x} + 1\right)$; \qquad $u(x) =e^{x}$; \qquad $v(x) = -e^{2x} + 1$; \qquad $n(x) = \left(e^{2x} + 1\right)^{2}$\\[1.0ex]
                    $f''(x) = 0$\\[1.0ex]
                    \begin{minipage}[t]{0.4\linewidth}
                        \vspace{-2.0ex}
                        \begin{align*}
                            v(x) &= 0 \\[2.0ex]
                            -e^{2x} + 1 &= 0 &&| \quad + e^{2x}\\[1.5ex]
                            e^{2x} &= 1 &&\\[1.5ex]
                            2x &= \ln\left(1\right) &&| \quad :2\\[1.5ex]
                            x_{1} &= 0 \ \text{\textit{(e)}}&&\\[-2.0ex]
                        \end{align*}
                    \end{minipage}
                    \hfill
                    \begin{minipage}[t]{0.55\linewidth}
                        \vspace{-2.0ex}
                        \begin{align*}
                            n(x) > 0
                        \end{align*}
                    \end{minipage}
                    \raisebox{-\height}{
                    \begin{tabularx}{0.95\linewidth}{c|X|X|X}
                        & $-\infty < x < 0$ & $x = 0$ & $0 < x < \infty$ \\ \hline
                        $u(x)$ & $+$ & $+$ & $+$ \\ \hline
                        $v(x)$ & $+$ & $0$ & $-$ \\ \hline
                        $z(x)$ & $+$ & $0$ & $-$ \\ \hline
                        $n(x)$ & $+$ & $+$ & $+$ \\ \hline
                        $f''(x)$ & $-$ & $0$ & $-$ \\ \hline
                        $G_{f}$ & l. g. & $W_{1}$ & r. g. \\
                    \end{tabularx}
                    }\\[2.0ex]
                    $\Rightarrow \qquad G_{f}$ ist linksgekrümmt in $\left] \ -\infty; 0 \ \right]$ \\[1.0ex]
                    $\Rightarrow \qquad G_{f}$ ist rechtsgekrümmt in $\left[ \ 0; +\infty \ \right[$\\[2.0ex]
                    $W_{1} \left( 0 \ / \ \dfrac{\pi}{4} \right)$\\[-2.0ex]
                    }
                \item Wertemenge: \\[1.0ex]{
                    Aus Grenzverhalten und Monotonie folgt $W_{f} = \left] \ 0; \dfrac{\pi}{2} \ \right[$
                }
                \item Zeichnung: \\[-2.0ex]{
                    \raisebox{-\height}{
                    \begin{tikzpicture}[
                            declare function={
                                f_f(\x) = rad(atan(exp(\x*\xScale)))*(1/\yScale);
                            }
                        ]

                        % Set variables
                        \pgfmathsetmacro{\axisPosWidth}{6.0};
                        \pgfmathsetmacro{\axisNegWidth}{-6.0};
                        \pgfmathsetmacro{\axisPosHeight}{4.0};
                        \pgfmathsetmacro{\axisNegHeight}{-0.5};
                        \pgfmathsetmacro{\xIncrement}{1.0};
                        \pgfmathsetmacro{\yIncrement}{1.0};
                        \pgfmathsetmacro{\xScale}{1.0};
                        \pgfmathsetmacro{\yScale}{0.5};
                        \pgfmathsetmacro{\gridxIncrement}{0.5};
                        \pgfmathsetmacro{\gridyIncrement}{0.5};

                        % Axis/grid limits
                        \pgfmathsetmacro{\xUpperLimitAxis}{\xIncrement * floor(\axisPosWidth / \xIncrement)}
                        \pgfmathsetmacro{\xLowerLimitAxis}{\xIncrement * ceil(\axisNegWidth / \xIncrement)}
                        \pgfmathsetmacro{\yUpperLimitAxis}{\yIncrement * floor(\axisPosHeight /\yIncrement)}
                        \pgfmathsetmacro{\yLowerLimitAxis}{\yIncrement * ceil(\axisNegHeight / \yIncrement)}

                        \pgfmathsetmacro{\xUpperLimitGrid}{\gridxIncrement * floor(\axisPosWidth / \gridxIncrement)}
                        \pgfmathsetmacro{\xLowerLimitGrid}{\gridxIncrement * ceil(\axisNegWidth / \gridxIncrement)}
                        \pgfmathsetmacro{\yUpperLimitGrid}{\gridyIncrement * floor(\axisPosHeight / \gridyIncrement)}
                        \pgfmathsetmacro{\yLowerLimitGrid}{\gridyIncrement * ceil(\axisNegHeight / \gridyIncrement)}

                        % Axis environment (PGFPlots)
                        \begin{axis}[
                            xmin=\xLowerLimitGrid, xmax=\xUpperLimitGrid,
                            ymin=\yLowerLimitGrid, ymax=\yUpperLimitGrid,
                            axis lines=middle,
                            axis line style={->, >=stealth},
                            grid=both,
                            xlabel={$x$},
                            ylabel={$y$},
                            xlabel style={anchor=west, at={(current axis.right of origin)}},
                            ylabel style={anchor=south, at={(current axis.above origin)}},
                            xtick={\xLowerLimitAxis,\xLowerLimitAxis+1,...,\xUpperLimitAxis},
                            ytick={\yLowerLimitAxis,\yLowerLimitAxis+1,...,\yUpperLimitAxis},
                            xticklabel={\pgfmathparse{\tick*\xScale}\pgfmathprintnumber{\pgfmathresult}},
                            yticklabel={\pgfmathparse{\tick*\yScale}\pgfmathprintnumber{\pgfmathresult}},
                            tick label style={font=\scriptsize},
                            x=1cm,
                            y=1cm,
                            minor tick num=1,
                            scale only axis,
                            grid=both,
                            restrict y to domain=\yLowerLimitGrid:\yUpperLimitGrid,
                        ]
                            
                        % Graph limits
                        \pgfmathsetmacro{\xLowerLimitGraph}{\xLowerLimitAxis}
                        \pgfmathsetmacro{\xUpperLimitGraph}{\xUpperLimitAxis}
                        
                        % G_f
                        \pgfmathsetmacro{\Gfax}{-3*(1/\xScale)}
                        \pgfmathsetmacro{\Gfay}{f_f(\Gfax)}
                        \addplot[thick, myb, samples=1000, domain=\xLowerLimitGraph:\xUpperLimitGraph]
                            {f_f(x)};
                        \node[anchor=south east] at (axis cs:{\Gfax},{\Gfay}) {\textcolor{myb}{$G_{f}$}};

                        % W_1
                        \pgfmathsetmacro{\Wax}{0}
                        \pgfmathsetmacro{\Way}{f_f(\Wax)}
                        \pgfmathsetmacro{\WaxLabel}{\Wax*\xScale}
                        \pgfmathsetmacro{\WayLabel}{\Way*\yScale}
                        \addplot[only marks, mark=x, mark size=3pt] coordinates {(\Wax,\Way)};
                        \node[anchor=east, xshift=0ex, yshift=0ex] at (axis cs:{\Wax},{\Way}) {$W_{1}(\num[output-decimal-marker = {.},round-mode=places, round-precision=2]{\WaxLabel}{},\ \num[output-decimal-marker = {.},round-mode=places,round-precision=2]{\WayLabel}{})$};

                        % h_1 für y1
                        \pgfmathsetmacro{\Gtax}{3*(1/\xScale)}
                        \pgfmathsetmacro{\Gtay}{0*(1/\yScale)}
                        \addplot[thick, myorange] coordinates {(\xLowerLimitGraph,\Gtay) (\xUpperLimitGraph,\Gtay)};
                        \node[anchor=south west] at (axis cs:{\Gtax},{\Gtay}) {\textcolor{myorange}{$G_{h_{1}}$}};

                        % h_2 für y2
                        \pgfmathsetmacro{\Gtbx}{-3*(1/\xScale)}
                        \pgfmathsetmacro{\Gtby}{pi/2*(1/\yScale)}
                        \addplot[thick, myorange] coordinates {(\xLowerLimitGraph,\Gtby) (\xUpperLimitGraph,\Gtby)};
                        \node[anchor=south west] at (axis cs:{\Gtbx},{\Gtby}) {\textcolor{myorange}{$G_{h_{2}}$}};

                        \end{axis}
                    \end{tikzpicture}
                    }\\[2.0ex]
                    }
            \end{itemize}
            }
        \item[h)] {
            $f(x) = \arctan\left(e^{x} - 1\right)$; \qquad $u(x) = e^{x} - 1$ \\[-2.0ex]
            \begin{itemize}[leftmargin=*]
                \item Definitionsmenge: \\[1.0ex]{
                    $f(x)$ definiert für $u(x)$\\[1.0ex]
                    $\Rightarrow \quad \uuline{D_{f} = \mathbb{R}}$;
                    }
                \item Symmetrie: \\[-2.0ex]{
                    \begin{itemize}
                        \item Achsensymmetrie zur y-Achse: \quad $f(-x) = f(x)$\\[-2.0ex]{
                            \begin{align*}
                                \arctan\left(e^{-x} - 1\right) &= \arctan\left(e^{x} - 1\right) \quad \text{\textit{(falsch)}} &&
                            \end{align*}
                            }
                        \item Punktsymmetrie zum Ursprung: \quad $f(-x) = -f(x)$\\[-2.0ex]{
                            \begin{align*}
                                \arctan\left(e^{-x} - 1\right) &= -\arctan\left(e^{x} - 1\right) && \\[1.5ex]
                                \arctan\left(e^{-x} - 1\right) &= \arctan\left(-e^{x} + 1\right) \quad \text{\textit{(falsch)}} &&
                            \end{align*}
                            }
                    \end{itemize}
                    $\Rightarrow$ \quad keine Symmetrie\\[-1.0ex]
                    }
                \item Grenzverhalten: \\[-3.0ex]{
                    \begin{align*}
                        \lim\limits_{x \to -\infty} \arctan\left(\overbrace{e^{x} - 1}^{\textcolor{myg}{-1^{+}}}\right) &= -{\dfrac{\pi}{4}}^{+}; \qquad \qquad \lim\limits_{x \to +\infty} \arctan\left(\overbrace{e^{x}-1}^{\textcolor{myg}{+\infty}}\right) &= + {\dfrac{\pi}{2}}^{-} &&\\[-3.0ex]
                    \end{align*}
                    }
                \item Asymptoten: \\[-3.0ex]{
                    \begin{itemize}
                        \item Waagerechte Asymptoten: \quad {
                            \vspace{-5.0ex}
                            \begin{align*}
                                h_{1}(x): \quad &y = -\dfrac{\pi}{4}\\[1.5ex]
                                h_{2}(x): \quad &y = \dfrac{\pi}{2}
                            \end{align*}
                            }
                    \end{itemize}
                    }
                \item Schnittpunkte mit den Koordinatenachsen: \\[-2.0ex]{
                    \begin{itemize}
                        \item $y$-Achse: \quad $f(0) = \arctan\left(e^{0}-1\right) = 0$; \qquad $P_{1} \left( 0 \ / \ 0 \right)$\\[-1.0ex]
                        \item $x$-Achse: \quad $f(x) = 0$ \\[-2.0ex]{
                            \begin{align*}
                                \arctan\left(e^{x} - 1\right) &= 0 &&\\[1.5ex]
                                e^{x} - 1&= 0 &&| \quad + 1\\[1.5ex]
                                e^{x} &= 1 &&\\[1.5ex]
                                x&= \ln\left(1\right)\\[1.5ex]
                                x_{1} &= 0 \ \text{\textit{(e)}}&&\\[-3.0ex]
                            \end{align*}
                            $N_{1} \left( 0 \ / \ 0 \right)$
                            }
                    \end{itemize}
                    }
                \item Extrempunkte: \\[-3.0ex]{
                    \begin{align*}
                        f'(x) &= \dfrac{1}{1 + \left(e^{x} - 1\right)^{2}} \cdot e^{x} \cdot 1&&\\[1.5ex]
                        &\left(= \dfrac{e^{x}}{1 + e^{2x} - 2e^{x} +1}\right)&& \\[1.5ex]
                        &\left(= \dfrac{e^{x}}{e^{2x} - 2e^{x} + 2}\right)&& \\[1.5ex]
                        &= \dfrac{e^{x}}{1 + \left(e^{x} - 1\right)^{2}} &&\\[-2.0ex]
                    \end{align*}
                    $z(x) = e^{x}$; \qquad $n(x) = 1 + \left(e^{x} - 1\right)^{2}$\\[1.0ex]
                    $f'(x) = 0 \qquad \Rightarrow \qquad z(x) > 0 \qquad \Rightarrow$ \quad keine Nullstellen\\[1.0ex]
                    $n(x) > 0 \qquad \Rightarrow \qquad f'(x) > 0$\\[1.0ex]
                    $\Rightarrow \quad G_{f}$ ist sms in $\mathbb{R}$\\[1.0ex]
                    $\Rightarrow \quad G_{f}$ hat keine Extrempunkte\\[-2.0ex]
                    }
                \item Wendepunkte: \\[-2.0ex]{
                    \begin{align*}
                        f''(x) &= \dfrac{\left[1 + \left(e^{x} - 1\right)\right] \cdot e^{x} \cdot 1 - e^{x} \cdot 2 \cdot \left(e^{x} - 1\right) \cdot e^{x} \cdot 1}{\left[1 + \left(e^{x} + 1\right)^{2}\right]^{2}}&&\\[1.5ex]
                        &= \dfrac{e^{x} \cdot \left[1 + \left(e^{x} - 1\right) -  2 \cdot \left(e^{x} - 1\right) \cdot e^{x}\right]}{\left[1 + \left(e^{x} + 1\right)^{2}\right]^{2}}&&\\[1.5ex]
                        &= \dfrac{e^{x} \cdot \left[1 + e^{x} - 1 -  2e^{2x} + 2e^{x}\right]}{\left[1 + \left(e^{x} + 1\right)^{2}\right]^{2}}&&\\[1.5ex]
                        &= \dfrac{e^{x} \cdot \left[-2e^{2x} + 3e^{x}\right]}{\left[1 + \left(e^{x} + 1\right)^{2}\right]^{2}}&&\\[1.5ex]
                        &= \dfrac{e^{2x} \cdot \left(-2e^{x} + 3\right)}{\left[1 + \left(e^{x} + 1\right)^{2}\right]^{2}}&&\\[-2.0ex]
                    \end{align*}
                    $z(x) = e^{2x} \cdot \left(-2e^{x} + 3\right)$; \qquad $u(x) =e^{2x}$; \qquad $v(x) = -2e^{x} + 3$; \\[1.0ex]
                    $n(x) = \left[1 + \left(e^{x} + 1\right)^{2}\right]^{2}$\\[1.0ex]
                    $f''(x) = 0$\\[1.0ex]
                    \begin{minipage}[t]{0.4\linewidth}
                        \vspace{-2.0ex}
                        \begin{align*}
                            v(x) &= 0 \\[2.0ex]
                            -2e^{x} + 3 &= 0 &&| \quad + 2e^{x}\\[1.5ex]
                            2e^{x} &= 3 &&| \quad :2\\[1.5ex]
                            e^{x} &= \SI{1.5}{} &&\\[1.5ex]
                            x_{1} &= \ln\left(\SI{1.5}{}\right) \ \text{\textit{(e)}} &&\\[-1.5ex]
                        \end{align*}
                    \end{minipage}
                    \hfill
                    \begin{minipage}[t]{0.55\linewidth}
                        \vspace{-2.0ex}
                        \begin{align*}
                            n(x) > 0
                        \end{align*}
                    \end{minipage}
                    \raisebox{-\height}{
                    \begin{tabularx}{0.95\linewidth}{c|X|X|X}
                        & $-\infty < x < \ln\left(\SI{1.5}{}\right)$ & $x = \ln\left(\SI{1.5}{}\right)$ & $\ln\left(\SI{1.5}{}\right) < x < \infty$ \\ \hline
                        $u(x)$ & $+$ & $+$ & $+$ \\ \hline
                        $v(x)$ & $+$ & $0$ & $-$ \\ \hline
                        $z(x)$ & $+$ & $0$ & $-$ \\ \hline
                        $n(x)$ & $+$ & $+$ & $+$ \\ \hline
                        $f''(x)$ & $-$ & $0$ & $-$ \\ \hline
                        $G_{f}$ & l. g. & $W_{1}$ & r. g. \\
                    \end{tabularx}
                    }\\[2.0ex]
                    $\Rightarrow \qquad G_{f}$ ist linksgekrümmt in $\left] \ -\infty; \ln\left(\SI{1.5}{}\right) \ \right]$ \\[1.0ex]
                    $\Rightarrow \qquad G_{f}$ ist rechtsgekrümmt in $\left[ \ \ln\left(\SI{1.5}{}\right); +\infty \ \right[$\\[2.0ex]
                    $W_{1} \left( \SI{0.41}{} \ / \ \SI{0.46}{} \right)$\\[-2.0ex]
                    }
                \item Wertemenge: \\[1.0ex]{
                    Aus Grenzverhalten und Monotonie folgt $W_{f} = \left] \ -\dfrac{\pi}{4}; \dfrac{\pi}{2} \ \right[$
                }
                \item Zeichnung: \\[-2.0ex]{
                    \raisebox{-\height}{
                    \begin{tikzpicture}[
                            declare function={
                                f_f(\x) = rad(atan(exp(\x*\xScale)-1))*(1/\yScale);
                            }
                        ]

                        % Set variables
                        \pgfmathsetmacro{\axisPosWidth}{6.0};
                        \pgfmathsetmacro{\axisNegWidth}{-6.0};
                        \pgfmathsetmacro{\axisPosHeight}{4.0};
                        \pgfmathsetmacro{\axisNegHeight}{-2.5};
                        \pgfmathsetmacro{\xIncrement}{1.0};
                        \pgfmathsetmacro{\yIncrement}{1.0};
                        \pgfmathsetmacro{\xScale}{1.0};
                        \pgfmathsetmacro{\yScale}{0.5};
                        \pgfmathsetmacro{\gridxIncrement}{0.5};
                        \pgfmathsetmacro{\gridyIncrement}{0.5};

                        % Axis/grid limits
                        \pgfmathsetmacro{\xUpperLimitAxis}{\xIncrement * floor(\axisPosWidth / \xIncrement)}
                        \pgfmathsetmacro{\xLowerLimitAxis}{\xIncrement * ceil(\axisNegWidth / \xIncrement)}
                        \pgfmathsetmacro{\yUpperLimitAxis}{\yIncrement * floor(\axisPosHeight /\yIncrement)}
                        \pgfmathsetmacro{\yLowerLimitAxis}{\yIncrement * ceil(\axisNegHeight / \yIncrement)}

                        \pgfmathsetmacro{\xUpperLimitGrid}{\gridxIncrement * floor(\axisPosWidth / \gridxIncrement)}
                        \pgfmathsetmacro{\xLowerLimitGrid}{\gridxIncrement * ceil(\axisNegWidth / \gridxIncrement)}
                        \pgfmathsetmacro{\yUpperLimitGrid}{\gridyIncrement * floor(\axisPosHeight / \gridyIncrement)}
                        \pgfmathsetmacro{\yLowerLimitGrid}{\gridyIncrement * ceil(\axisNegHeight / \gridyIncrement)}

                        % Axis environment (PGFPlots)
                        \begin{axis}[
                            xmin=\xLowerLimitGrid, xmax=\xUpperLimitGrid,
                            ymin=\yLowerLimitGrid, ymax=\yUpperLimitGrid,
                            axis lines=middle,
                            axis line style={->, >=stealth},
                            grid=both,
                            xlabel={$x$},
                            ylabel={$y$},
                            xlabel style={anchor=west, at={(current axis.right of origin)}},
                            ylabel style={anchor=south, at={(current axis.above origin)}},
                            xtick={\xLowerLimitAxis,\xLowerLimitAxis+1,...,\xUpperLimitAxis},
                            ytick={\yLowerLimitAxis,\yLowerLimitAxis+1,...,\yUpperLimitAxis},
                            xticklabel={\pgfmathparse{\tick*\xScale}\pgfmathprintnumber{\pgfmathresult}},
                            yticklabel={\pgfmathparse{\tick*\yScale}\pgfmathprintnumber{\pgfmathresult}},
                            tick label style={font=\scriptsize},
                            x=1cm,
                            y=1cm,
                            minor tick num=1,
                            scale only axis,
                            grid=both,
                            restrict y to domain=\yLowerLimitGrid:\yUpperLimitGrid,
                        ]
                            
                        % Graph limits
                        \pgfmathsetmacro{\xLowerLimitGraph}{\xLowerLimitAxis}
                        \pgfmathsetmacro{\xUpperLimitGraph}{\xUpperLimitAxis}
                        
                        % G_f
                        \pgfmathsetmacro{\Gfax}{-3*(1/\xScale)}
                        \pgfmathsetmacro{\Gfay}{f_f(\Gfax)}
                        \addplot[thick, myb, samples=1000, domain=\xLowerLimitGraph:\xUpperLimitGraph]
                            {f_f(x)};
                        \node[anchor=south east] at (axis cs:{\Gfax},{\Gfay}) {\textcolor{myb}{$G_{f}$}};

                        % N_1
                        \pgfmathsetmacro{\Nax}{0}
                        \pgfmathsetmacro{\Nay}{f_f(\Nax)}
                        \pgfmathsetmacro{\NaxLabel}{\Nax*\xScale}
                        \pgfmathsetmacro{\NayLabel}{\Nay*\yScale}
                        \addplot[only marks, mark=x, mark size=3pt] coordinates {(\Nax,\Nay)};
                        \node[anchor=south east, xshift=0ex, yshift=0ex] at (axis cs:{\Nax},{\Nay}) {$N_{1}(\num[output-decimal-marker = {.},round-mode=places, round-precision=2]{\NaxLabel}{},\ \num[output-decimal-marker = {.},round-mode=places,round-precision=2]{\NayLabel}{})$};

                        % W_1
                        \pgfmathsetmacro{\Wax}{ln(1.5)}
                        \pgfmathsetmacro{\Way}{f_f(\Wax)}
                        \pgfmathsetmacro{\WaxLabel}{\Wax*\xScale}
                        \pgfmathsetmacro{\WayLabel}{\Way*\yScale}
                        \addplot[only marks, mark=x, mark size=3pt] coordinates {(\Wax,\Way)};
                        \node[anchor=north west, xshift=0ex, yshift=0ex] at (axis cs:{\Wax},{\Way}) {$W_{1}(\num[output-decimal-marker = {.},round-mode=places, round-precision=2]{\WaxLabel}{},\ \num[output-decimal-marker = {.},round-mode=places,round-precision=2]{\WayLabel}{})$};

                        % h_1 für y1
                        \pgfmathsetmacro{\Gtax}{3*(1/\xScale)}
                        \pgfmathsetmacro{\Gtay}{-pi/4*(1/\yScale)}
                        \addplot[thick, myorange] coordinates {(\xLowerLimitGraph,\Gtay) (\xUpperLimitGraph,\Gtay)};
                        \node[anchor=south west] at (axis cs:{\Gtax},{\Gtay}) {\textcolor{myorange}{$G_{h_{1}}$}};

                        % h_2 für y2
                        \pgfmathsetmacro{\Gtbx}{-3*(1/\xScale)}
                        \pgfmathsetmacro{\Gtby}{pi/2*(1/\yScale)}
                        \addplot[thick, myorange] coordinates {(\xLowerLimitGraph,\Gtby) (\xUpperLimitGraph,\Gtby)};
                        \node[anchor=south west] at (axis cs:{\Gtbx},{\Gtby}) {\textcolor{myorange}{$G_{h_{2}}$}};

                        \end{axis}
                    \end{tikzpicture}
                    }
                    }
            \end{itemize}
            }
    \end{enumerate}
\end{enumerate}
\end{Exercise}